\section{Conclusion}
\label{sec:conclusion}

We have argued that the structural failure of direct-generation silicon sampling---miscalibration, prompt sensitivity, and poor transfer---stems from forcing an LLM to commit to a discrete answer token, which discards the rich latent structure of how a respondent would react. Drawing an explicit analogy to the use of LLMs as representation generators in spatio-temporal learning, we proposed to treat the LLM as a latent respondent encoder rather than an answer generator. Our method, ResponseVec, fuses a demographic profile, target-relevant answer history, and the current item with its options into a literature-grounded survey prompt, extracts a latent reaction vector with a generative LLM embedder that preserves output semantics, and predicts the true choice with a lightweight, calibratable, option-aware decoder. We further introduced a reusable RespondentVec variant that amortizes compute across items, and we formalized the three-way trade-off among task adaptivity, representation reusability, and compute cost. Across three transfer regimes---unseen questions, cross-domain items, and out-of-distribution demographic intersections---we designed experiments to test whether the representation-based factorization is more accurate, better calibrated, more robust, and more transferable than seven baseline families spanning direct generation, classical psychometric models, item-conditional collaborative filtering, quantized fine-tuning, generic text vectors, raw hidden states, and a discriminative-encoder ablation. Our anticipated results, organized around four questions and three predictions, suggest that the representation-based approach traces a Pareto frontier that uniform direct-generation methods cannot reach, and that the frontier shifts systematically with the transfer regime in a way that yields operational guidance for practitioners. The principal contributions are conceptual (the predictor-to-representation-generator transposition to the survey domain), methodological (the option-aware calibratable decoder and the use of a generative embedder for latent reaction), empirical (the first three-way transfer evaluation against nine baselines, filling the future-wave gap), and analytical (the characterization of the ResponseVec/RespondentVec boundary). Future work should extend the evaluation to longitudinal panel surveys with genuine temporal transfer, audit the contamination exposure of the generative embedder, and explore interactions among the design choices we isolated. More broadly, we hope this work reframes the role of LLMs in survey research---from answer generators that compete with human respondents to representation generators that augment them with calibrated, transferable, and cost-aware imputations.
